\section{Discussion}
\label{sec:discussion}

\subsection{What the controlled evidence says}

% CLAIM: C03,C04,C05,C06,C07,C08
The analytic and development results are most informative when read together.
Reliable analytic velocity sharply reduces position error and grid lag for the
three smooth references, establishing that target-state information can matter
under a controlled follower and timing convention.  It does not establish
that a derivative estimated from arbitrary positions will be sufficiently
accurate, timely, or predictive.  The development trace supplies that missing
warning: all tested finite-difference-derived pipelines are worse than the
position-only baseline.  Because the fixed projection acts on the PVA targets
but not the P or PV targets, the experiment cannot causally allocate the
degradation between differentiation and projection.  It nevertheless rejects
the operational proposition that these complete tested pipelines improve this
trace.

The PV/PVA equality is narrower still.  The two conditions reach the same
requested position and velocity endpoint, and their different requested
accelerations do not change the recorded position endpoint in the tested
ordinary state-to-state solves.  This mechanism is not evidence that
acceleration is generally irrelevant, nor is it a comparison of acceleration
estimators, internal profiles, or control effort.  The smooth Phase~A
references do not represent acceleration-active discontinuities or regimes in
which the terminal acceleration changes a binding constraint.  Noisy or
inadmissible acceleration estimates may be harmful because they alter the
requested state and may trigger projection, but that explanation remains a
hypothesis outside the observed fixed pipeline.

The noncausal next-cycle oracle clarifies a different issue.  Near-zero error
and zero grid lag occur when the target belongs to the output time and is
one-step reachable.  That sanity control supports the time-index contract; it
is not a deployable predictor or a bound on every causal method.  In practice,
estimation and future-reference generation should remain separate interfaces
so that delay, horizon, model assumptions, and target time can be inspected
rather than absorbed into a derivative label.

\subsection{Execution semantics before performance ranking}

% CLAIM: C01,C02,C09,C11,C12
A continuously updated endpoint can be processed by a state-to-state OTG, but
that does not make the operation semantically identical to following a moving
reference.  The reference generator decides what future state is desired; the
governor decides what adjacent endpoint is executable; the follower decides
which continuous profile reaches the committed endpoint.  These layers can be
implemented together, but their responsibilities remain distinct for analysis
and logging.

The direct construction makes the adjacent-state contract unusually visible:
the selected jerk generates the endpoint by exact integration.  A second
trajectory-shaping stage can be useful, but it changes the executed profile
and must be audited on its own terms.  In particular, a \Ruckig{} prefix can
switch jerk within one control period.  Replacing it by the endpoint
acceleration difference loses the segment sequence and cannot establish
internal-jerk legality.  When exact segments are exposed, continuous velocity,
acceleration, and jerk extrema can be checked analytically; when only sampled
states are available, internal jerk must remain marked unavailable.

Method identity follows the command.  An unshielded native prefix is
\OrdinaryRuckig{}.  A shield that may replace it defines a
\ViabilityShieldedRuckig{} condition.  Executing the governor jerk directly is
\DirectConstantJerkExecution{}.  An algorithm-changing fallback is not merely
a status code: it changes which controller generated the command and can
invalidate a nominal comparison.  This identity rule applies independently of
whether the replacement is safer or performs better.

\subsection{Evidence correction and method identity}
\label{sec:evidence-correction}

% CLAIM: C12,C13,E01,N03
The frozen \(\VThreeExploratoryConfoundedImprovement\)\% value is retained only
as an exposed exploratory regression: the historical \texttt{predicted\_p}
condition used an algorithm-changing fallback on
\(\VThreeConfoundedBaselineFallbackCount/
\VThreeConfoundedBaselineCycleCount\) cycles
(\(\VThreeConfoundedBaselineFallbackRate\%\)), so the comparison is neither
\OrdinaryRuckig{} nor a
same-follower P/PVA ablation; its confirmatory classification is withdrawn,
the v3 artifact bytes remain unchanged, and v3 was not rerun.  This correction
does not change the separately recorded direct constant-jerk audits.  It does
show why immutable artifacts, executed-command identity, and denominators are
needed together: checksums preserve what happened, while scientific
classification determines what that evidence can support.

\subsection{Scope and required evidence}

% CLAIM: C10,N01,N02,N03
The frozen direct condition supports a command-model audit, not a general
safety probability.  Its zero counts are conditional on a synthetic
triple-integrator setting, known limits, viable selected current states, exact
execution of the chosen jerk, and the registered numerical tolerance.  The
runtime measurements apply to one software/hardware environment and a
separate post-warm-up population.  They provide no evidence of real-time behavior on
another processor, within a robot control stack, or under communication and
measurement load.

The one development CSV is position-only, fixed to one row per
\(\ControlPeriodMS\) ms, and not an independent test set.  Its stored
timestamps are outside the reported replay semantics.  Nothing here validates
irregular sampling, source-clock reconstruction, resampling, network jitter,
or real motion derivatives.  Likewise, the kinematic model omits torque,
actuator bandwidth, coupling, collision, contact, thermal constraints,
disturbance rejection, and measured closed-loop plant behavior.

The decisive next experiment is not another mixed baseline.  A fresh locked
study must hold the executable governor and follower fixed while varying only
P, PV, and PVA information, use new experiment identities and seeds, expose
acceptance and fallback rates, and pre-specify its estimators and target
times.  An independent position stream should then test estimator and
prediction robustness, followed by hardware or hardware-in-the-loop work that
adds dynamics, measurement, computation, and safety-system boundaries.  Until
those steps, the durable result is the timing/feasibility formulation and its
scoped command-execution evidence, not a performance ranking.
